% Part: normal-modal-logic
% Chapter: frame-correspondence
% Section: first-order-definability

\documentclass[../../../include/open-logic-section]{subfiles}

\begin{document}

\olfileid[zh]{nml}{frd}{fol}

\olsection{一阶可定义性}

我们已经看到，框架的许多可及关系性质可以由模态!!{formula}定义。例如，框架的对称性可以由!!{formula}~\Ax{B}，$p \lif \Box\Diamond
p$ 定义。我们迄今遇到的各种条件都可以由只含一个二元!!{predicate}的!!a{language}中的一阶!!{formula}表达。例如，对称性由
$\lforall[x][\lforall[y][(\Atom{Q}{x,y} \lif \Atom{Q}{y, x})]]$ 定义，意思是：当且仅当 $R$ 对称时，一个满足 $\Domain{M} =
W$ 与 $\Assign{Q}{M} = R$ 的一阶!!{structure}~$\Struct{M}$ 满足前面的公式。这便引出了如下定义：

\begin{defn}
  一个框架类~$\mClass{F}$ 若存在一个一阶语言中含单个二元!!{predicate}~$Q$ 的!!a{sentence}~$!A$，使得 $\mModel{F} =
  \tuple{W, R} \in \mClass{F}$ 当且仅当在满足 $\Domain{M} = W$ 与 $\Assign{Q}{M}
  = R$ 的一阶!!{structure}~$\Struct{M}$ 中，$\Sat{M}{!A}$
  ，则称它是\emph{一阶可定义的}。
\end{defn}

事实证明，迄今考虑的这些性质以及定义它们的模态!!{formula}都是特殊情形。并非每个!!{formula}都能定义一个一阶可定义的框架类，也并非每个一阶可定义的框架类都能由某个模态公式定义。

对第一种情形给出反例的是 L\"ob 公式：
\begin{equation}
\Box(\Box p \lif p) \lif \Box p. \tag{\Ax{W}}
\end{equation}
\Ax{W} 定义传递且反良基的框架类。一个关系是良基的，如果不存在无穷序列
$w_1$, $w_2$, \dots{} 使得 $Rw_2w_1$, $Rw_3w_2$, \dots。例如，$\Nat$~上的关系 $<$ 是良基的，而 $\Int$~上的关系 $<$ 则不是。一个关系是反良基的，当且仅当其逆关系是良基的。因此，反良基关系就是那些不存在无穷序列 $w_1$, $w_2$, \dots{}
使得 $Rw_1w_2$, $Rw_2w_3$, \dots 的关系。

然而，并不存在一阶!!{formula}来定义传递的
反良基关系。因为设 $\Sat{M}{!F}$ 当且仅当 $R =
\Assign{Q}{M}$ 是传递且反良基的。令 $!A_n$ 为!!{formula}
\[
(\Atom{Q}{a_1,a_2} \land \dots \land
    \Atom{Q}{a_{n-1},a_{n}})
\]
现在考虑!!{formula}的集合
\[
\Gamma = \{!F, !A_1, !A_2, \dots\}.
\]
$\Gamma$ 的每个有穷子集都是可满足的：令 $k$ 为使得 $!A_k$ 在该子集中出现的最大下标，$\Domain{M_k} = \{1, \dots, k\}$，
$\Assign{a_i}{M_k} = i$，且 $\Assign{Q}{M_k} = <$。由于 $\{1,
\dots, k\}$ 上的 $<$ 是传递且反良基的，
故由构造知对所有 $i \le k$ 有 $\Sat{M_k}{!F}$ 且 $\Sat{M_k}{!A_i}$。
由一阶逻辑的紧致性定理，$\Gamma$ 在某个!!{structure}~$\Struct{M}$ 中是可满足的。由假设，由于
$\Sat{M}{!F}$，关系 $\Assign{Q}{M}$ 是反
良基的。但显然，$\Assign{a_1}{M}$, $\Assign{a_2}{M}$,
\dots{} 会构成一个被反良基性所排除的那种无穷序列。

对第二种论断给出反例的是全域性这个性质：对所有 $u$ 与 $v$，都有 $Ruv$。全域框架可由公式
$\lforall[x][\lforall[y][\Atom{Q}{x,y}]]$ 一阶地定义。然而，没有任何模态公式恰好只在全域框架中有效。这是一个独立有趣的结论的后承：在全域框架中有效的公式恰好与在自反、对称且传递的框架中有效的公式相同。存在并非全域的自反、对称且传递的框架，因此每个
在所有全域框架中都有效的!!{formula}也在某些非全域
框架中有效。

\end{document}
